\scanpage{30}
\[
\deleted{\boxed{\text{\textcircled{F}}\otimes G}}
\qquad
L.(F)\otimes L.(G)=(F\underline{\underline{\otimes}}G)
\]
\margin{\figure{\begin{tikzpicture}
\node[anchor=west] at (.2,4) {$n=2\nu$};
\draw (0,4.4)--(-.1,2.4);
\draw[->] (-.12,1.9) .. controls (-.2,1.5) and (-.25,.3) .. (-.3,0);
\end{tikzpicture}}}
\diagram{\begin{tikzcd}[column sep=large,row sep=large]
H^{n+1}(U^{n+1})(1)
\arrow[r,"\mathrm{inj}","\partial"']
\arrow[rr,bend left=35,"{\deleted{\unsure{bij}}\ \mathrm{bij}\ \uline{exc.\ 2},\ \ell^{(2\nu+1)}\rightsquigarrow2\ell^{(2\nu)}}"]
&H^n(Q^n)\arrow[r,"\mathrm{surj}"]&H^n(U^n)\\
\ell^{2\nu+1}\arrow[r,rightsquigarrow]&\lambda_\nu\\
&{[\lambda_\nu,\xi^\nu]}\arrow[r,rightsquigarrow]&{[2\ell_{\deleted{\unsure{2\nu}}}^{2\nu},0]}\\
&\varphi_\nu=\frac12(\lambda_\nu+\xi^\nu)\arrow[r,rightsquigarrow]&\ell_{\deleted{\unsure{2\nu}}}^{2\nu}
\end{tikzcd}}
\margin{$0\text{---}$}
\[
\ell_{\deleted{\unsure{2\nu+1}}}^{2\nu+1}
\]
\[
\deleted{H^n(Q^n)\longleftarrow H^n(U^n)}
\]
\diagram{\begin{tikzcd}[column sep=large,row sep=large]
&\lambda_\nu&\ell_{2\nu}\arrow[l,rightsquigarrow]\\
H_c^{n+1}(U^{n+1})(+1)&H^n(Q^n)\arrow[l,"\mathrm{surj}","\partial"']&H_c^n(U^n)\arrow[l,"\mathrm{inj}"]\\
{[2\ell_{2\nu+1},0]}&{[\lambda_\nu,\xi^\nu]}\arrow[l,rightsquigarrow]\\
\ell_{2\nu+1}&\varphi_\nu\arrow[l,rightsquigarrow]\\
{}&&{}\arrow[ll,bend left=30]\\
2\ell_{2\nu+1}&&\ell_{2\nu}\arrow[ll,rightsquigarrow]
\end{tikzcd}}
\[
Q^{n+1}-Q^n=U^{n+1}
\]

\uline{N.B.}\quad \uline{On \unsure{interprète}} $\unsure{\varphi_\nu}$ \uline{directement par les sous-espaces isotropes maximaux.}

les deux familles de tels sous-espaces \unsure{correspondent|corresp.}: les classes $\alpha_\nu,\beta_\nu$. On a
\[
\alpha_\nu+\beta_\nu=\xi^\nu,\qquad
\alpha_\nu-\beta_\nu=\lambda_\nu,\qquad
\varphi_\nu=\alpha_\nu\text{---}
\]
\mspage{29}
